\section{Introduction}

\subsection{Motivation}

Sparse attention has seen rapid progress in 2024--2025. Notable recent 
advances include:

\begin{itemize}
\item \textbf{XAttention} (Xu et al., ICML 2025): training-free block scoring via
  antidiagonal correlations.
\item \textbf{SeerAttention} (Gao et al., NeurIPS 2025): learnable per-head block
  sparsity.
\item \textbf{DuSA} (Cao et al., NeurIPS 2025): dual-stage intra-block dense +
  inter-block sparse.
\item \textbf{AnchorAttention} (Zhu et al., 2025): stripe-based fine-grained sparsity.
\item \textbf{Spectraformer} (Nguyen et al., 2025): RFF-based global rectification.
\item \textbf{Softpick} (Zuhri et al., 2025): rectified softmax with natural ~47\%
  sparsity.

\end{itemize}
\textbf{These techniques are complementary but rarely combined.} NEXUS is our 
attempt to unify them in a single attention module.

\subsection{Problem Statement}

Given a pretrained transformer $\mathcal{M}$ with $H$ heads and hidden 
dimension $d$, we seek an attention mechanism $\mathcal{A}'$ such that:

\[
\text{FLOPs}(\mathcal{A}') \ll \mathcal{O}(N^2), \quad
\mathcal{A}' \approx \mathcal{A}
\tag{1.1}
\]

subject to:

\begin{enumerate}
\item \textbf{Modularity:} Each sparse mechanism can be enabled/disabled.
\item \textbf{Multi-architecture:} DistilBERT, RoBERTa, GPT-2.
\item \textbf{Pretrained compatibility:} Reuse pretrained Q/K/V/O weights.
\item \textbf{Causal correctness:} Preserve causality for decoders.
\item \textbf{Training stability:} No collapses.

\end{enumerate}
\subsection{Our Approach}

NEXUS combines six sparse-attention mechanisms in a unified pipeline:

\subsubsection{Block Partitioning}

The sequence is partitioned into $M = \lceil N / B_s \rceil$ blocks of size 
$B_s$ (default 16).

\subsubsection{XAttention Block Scoring}

Sub-sampled queries and keys compute coarse block-to-block scores:

\[
S_{b_i, b_j} = \frac{1}{|B_i| |B_j|} \sum_{i \in B_i, j \in B_j} \frac{q_i \cdot k_j}{\sqrt{d_h}}
\tag{1.2}
\]

Top-$k$ key blocks are retained per query block.

\subsubsection{SeerAttention Learnable Gating}

A learnable bottleneck network predicts per-head block budgets:

\[
g_h = \sigma(W_q \bar{Q}_h \cdot W_k \bar{K}_h^T)
\tag{1.3}
\]

The gate overrides/adds to the XAttention mask.

\subsubsection{DuSA Dual-Stage}

Intra-block (diagonal) is always dense. Inter-block uses top-$k$.

\subsubsection{AnchorAttention Stripe}

For $N > 256$, a stripe band mask restricts fine-grained within-block 
attention to adjacent stripe groups.

\subsubsection{Softpick}

A rectified softmax:

\[
\text{softpick}(x)_i = \frac{\text{ReLU}(e^{x_i/T} - 1)}{\sum_j \text{ReLU}(e^{x_j/T} - 1)}
\tag{1.4}
\]

\subsubsection{Spectraformer RFF Rectification}

Periodically, a global linear attention with RFF features is applied.

\subsection{Contributions}

\begin{itemize}
\item \textbf{A unified framework} for combining six recent sparse attention
  mechanisms in a single module.
\item \textbf{Modular design} with per-component enable/disable flags.
\item \textbf{Multi-architecture support} for flat, nested, and causal transformers.
\item \textbf{Honest complexity analysis.} We report which components are active under
  which conditions.
\item \textbf{Reproducible ablation.} Each component is ablated independently.

\end{itemize}
\subsection{Paper Organization}

Section 2 reviews related work. Section 3 provides background on each of 
the six components. Section 4 presents NEXUS. Section 5 describes 
integration. Section 6 details setup. Section 7 reports results. Section 8 
analyzes per-component contributions. Section 9 discusses trade-offs. 
Section 10 lists limitations. Section 11 concludes.



\subsection{Notation}

\begin{lstlisting}

---
\end{lstlisting}